\documentclass{article}

\usepackage[final]{neurips_2019}

\usepackage[utf8]{inputenc}
\usepackage[T1]{fontenc}
\usepackage{hyperref}
\usepackage{url}
\usepackage{booktabs}
\usepackage{amsfonts}
\usepackage{amsmath}
\usepackage{nicefrac}
\usepackage{microtype}
\usepackage{graphicx}
\usepackage{xcolor}

\graphicspath{{../figures/}}

\newcommand{\note}[1]{\textcolor{blue}{{#1}}}
\newcommand{\ET}{E_T}
\newcommand{\ER}{E_R}

\title{
  \textbf{Test-Time Confidence Prediction for Spacecraft Pose
Estimation Under Domain Shift} \\
  \vspace{1em}
  \small{\normalfont Stanford CS229 Project}
}

\author{
  Anjali Sreenivas \\
  Department of Computer Science \\
  Stanford University \\
  \texttt{anjalisr@stanford.edu} \\
  \And
  Lundeen Cahilly \\
  Department of Computer Science \\
  Stanford University \\
  \texttt{lcahilly@stanford.edu} \\
}

\begin{document}

\maketitle

\begin{abstract}
Learning-based spacecraft pose estimators are trained on synthetic imagery but
deployed on real images, where accuracy degrades sharply and unpredictably. A
downstream navigation filter must know \emph{which} pose estimates to trust, yet
modern estimators provide no such signal. We ask whether the reliability of a
state-of-the-art estimator (SPNv2) can be predicted at test time using only
features read off a single forward pass---and trained without any labels from the
deployment domain. Treating reliability as binary failure prediction against
absolute physical thresholds (translation $\ET>0.10$\,m, rotation $\ER>3.0^\circ$),
we train a class-weighted logistic regression on synthetic labels and evaluate on
two hardware-in-the-loop (HIL) domains. The predictor matches a HIL-trained oracle
(translation AUC $0.84$/$0.80$ vs.\ $0.85$/$0.80$) with \emph{zero} target labels,
showing that the feature$\rightarrow$failure map transfers across the domain gap
even though the pose estimator itself does not. As a frozen accept/reject gate it
isolates real physical error: rejected images carry $\sim$3--$4.5\times$ the
translation error of accepted ones. Importance weighting for covariate shift is
inert here---a free insurance policy that would activate under more severe shift.
\end{abstract}

\section{Introduction}

Autonomous pose estimation of noncooperative spacecraft is essential for
safety-critical missions such as on-orbit servicing, debris removal, and satellite
refueling \citep{pauly_survey_2023}. Accurate 6D pose estimation from monocular
imagery enables safe autonomous rendezvous, proximity operations, and docking,
where even small errors can lead to collision or mission failure
\citep{park_speed_2022, pasqualetto_cassinis_review_2019}. These pose estimates feed
downstream navigation filters such as the Unscented Kalman Filter (UKF), which rely
on \emph{per-measurement uncertainty} to weight observations and track a
spacecraft's state. Without reliable uncertainty, the filter cannot identify or
compensate for poor measurements, degrading navigation performance
\citep{park_robust_2024}.

This challenge is compounded by domain shift. Spaceborne ML models must be trained
on synthetic imagery (the only scalable source of labeled pose data), but they are
deployed on real images exhibiting harsh illumination, specular reflections, and
deep shadows that synthetic renderers struggle to reproduce
\citep{park_speed_2022, perez-villar_spacecraft_2024}. The Spacecraft Pose Network
v2 (SPNv2) \citep{park_robust_2024} is a state-of-the-art estimator that narrows
the synthetic-to-real gap via multi-task learning and online domain refinement, but
it explicitly identifies uncertainty quantification as unsolved, noting that
refinement worsens predictions on a substantial fraction of individual images even
when average performance improves.

\textbf{Input/output.} The input to our algorithm is a single monocular grayscale
image of the target spacecraft, summarized by a 13-dimensional feature vector read
directly from SPNv2's forward pass. We then use \emph{logistic regression}
(optionally importance-weighted for covariate shift) to output a predicted
reliability score: a calibrated failure probability for the translation and
rotation components of SPNv2's pose estimate, and a corresponding binary
accept/reject decision. The method modifies neither SPNv2 nor its weights and
requires no ground-truth pose labels from the deployment domain.

\section{Related Work}

\textbf{Spacecraft pose estimation and the domain gap.} Learning-based monocular
pose estimation has become the dominant approach for noncooperative targets
\citep{pauly_survey_2023, pasqualetto_cassinis_review_2019}. The KRN
\citep{park_towards_2019} and SPN families established the task on the SPEED/SPEED+
benchmarks \citep{park_speed_2022}; SPNv2 \citep{park_robust_2024} is the current
state of the art, using multi-task heads and online domain refinement, and
\citet{perez-villar_spacecraft_2024} push domain adaptation via inter-model
consensus. All of this work targets \emph{improving} the estimator; none predicts
\emph{per-image reliability} at test time, which is our focus.

\textbf{Uncertainty quantification and calibration.} Standard deep-learning
uncertainty tools include temperature/Platt scaling for calibration
\citep{guo_calibration_2017}, MC-dropout \citep{gal_dropout_2016}, and deep
ensembles \citep{lakshminarayanan_ensembles_2017}. The latter two require multiple
stochastic forward passes or retraining---prohibitive for a fixed, deployed
flight model---so we instead exploit single-pass, model-internal confidence
read-outs, in the spirit of using a model's own outputs as a misclassification/OOD
signal \citep{hendrycks_baseline_2017}.

\textbf{Covariate-shift adaptation.} When train and test inputs differ in
distribution but the conditional $p(y\mid x)$ is stable, importance weighting by
the density ratio $p_\text{test}(x)/p_\text{train}(x)$ yields consistent estimates
\citep{shimodaira_2000, sugiyama_2007}. We adopt the standard
domain-classifier estimator of this ratio and test whether it helps under the real
synthetic$\rightarrow$HIL shift.

\section{Dataset and Features}

\textbf{Dataset.} We use SPEED+ \citep{park_speed_2022}, from the Stanford Space
Rendezvous Laboratory, comprising 59{,}960 synthetic images and 9{,}531 HIL images
captured at Stanford's TRON testbed under calibrated spaceborne illumination, split
into diffuse Earth-albedo (\emph{lightbox}) and direct-sunlight (\emph{sunlamp})
conditions. HIL ground-truth poses are accurate to mm/sub-degree. We run SPNv2
inference and analyze 11{,}994 synthetic-validation, 6{,}737 lightbox, and 2{,}787
sunlamp images. Synthetic is the only labeled \emph{training} source; HIL labels are
used \emph{only} for evaluation (and to fit the oracle upper bound).

\textbf{Failure labels.} Following SPEED+ we measure translation error $\ET$ (m) and
rotation error $\ER$ (deg) separately and treat them as two independent failure
axes. We define failure with \emph{absolute, domain-independent} physical
thresholds---``a fail is a fail''---namely $\ET>0.10$\,m and $\ER>3.0^\circ$, rather
than domain-relative percentiles.

\textbf{The domain gap is severe and heavy-tailed (Fig.~\ref{fig:errdist}).} Mean
$\ET$ climbs $0.054\rightarrow0.175\rightarrow0.225$\,m and the translation failure
rate $11\%\rightarrow41\%\rightarrow63\%$ across synthetic$\rightarrow$lightbox
$\rightarrow$sunlamp. Rotation is worse: mean $\ER$ rises
$0.99^\circ\rightarrow6.4^\circ\rightarrow10.9^\circ$ with $p99$ near
$150$--$167^\circ$---the literal $180^\circ$ symmetry flips of the spacecraft body.

\textbf{Features.} Rather than image pixel statistics (which we found sign-flip
across the gap and fail to transfer), we read 13 features directly off SPNv2's
forward pass, so they live in \emph{pose/confidence space}, tied to the prediction
rather than the image's appearance: cross-head pose disagreement
(\texttt{disagree\_t\_m}, \texttt{disagree\_t\_norm}, \texttt{disagree\_R\_deg}),
segmentation and heatmap entropy (\texttt{seg\_entropy},
\texttt{hm\_entropy\_mean}), heatmap peak statistics (\texttt{hm\_peak\_mean},
\texttt{hm\_peak\_min}, \texttt{hm\_nconf}), detection confidence
(\texttt{effi\_cls}), and geometric covariates (\texttt{target\_distance},
\texttt{bbox\_area}, \texttt{bbox\_aspect}). \texttt{reject} has zero variance on
synthetic and is dropped. Remaining features span $\sim$7 orders of magnitude, so we
standardize each to $\mu{=}0,\sigma{=}1$ using synthetic-split statistics only (the
deployment-time information set).

\begin{figure}[t]
  \centering
  \includegraphics[width=0.92\linewidth]{pose_error_distribution.pdf}
  \caption{SPNv2 per-image error distributions across the three SPEED+ splits.
  Error grows and becomes heavy-tailed from synthetic (train) to the lightbox and
  sunlamp HIL domains, motivating per-image reliability prediction.}
  \label{fig:errdist}
\end{figure}

\section{Methods}

\textbf{Supervised confidence predictor.} For each axis we fit a logistic
regression $P(y{=}1\mid x)=\sigma(w^\top x + b)$ on synthetic features with binary
failure labels $y=\mathbf{1}[\text{error}>\tau]$, minimizing class-weighted
cross-entropy with $\ell_2$ regularization. Classes are balanced by inverse
frequency to counter the low synthetic failure rate ($11\%$/$2\%$ for
$\ET$/$\ER$). The fitted model is applied unchanged to HIL features. We compare
against an \textbf{oracle}---the same model \emph{trained on HIL labels}---to
upper-bound achievable performance and isolate the cost of the domain gap.

\textbf{Importance weighting (covariate-shift correction).} If only $p(x)$ shifts
while $p(y\mid x)$ is stable, reweighting the synthetic training loss by the density
ratio $w(x)=p_\text{HIL}(x)/p_\text{syn}(x)$ is consistent
\citep{shimodaira_2000}. We estimate $w$ with a domain classifier $d(x)\approx
P(\text{HIL}\mid x)$ (logistic regression, synthetic$=0$/HIL$=1$) via
$w(x)=d(x)/(1{-}d(x))$, clip extremes, and refit the failure model with these
weights. We report the effective sample size $\mathrm{ESS}=(\sum_i
w_i)^2/\sum_i w_i^2$ as a measure of how aggressively weighting reshapes the
training set.

\textbf{Baselines.} (i) \emph{Random}: predict failure at the domain's empirical
rate (AUC $=0.5$). (ii) \emph{Disagreement threshold}: use the raw physical-proxy
feature value as the score---\texttt{disagree\_t\_m} for translation,
\texttt{disagree\_R\_deg} for rotation---with no classifier. This is the cheap
physical baseline our learned model must beat. We skipped MC-dropout/ensembles
\citep{gal_dropout_2016, lakshminarayanan_ensembles_2017} as they require
re-running or retraining the frozen SPNv2; the single-pass confidence read-outs in
our feature set serve as the model's-own-confidence stand-in.

\textbf{Metrics.} Our primary metric is threshold-free \textbf{AUC-ROC}
(discrimination / ranking quality), reported per domain $\times$ axis. For
deployment we also freeze a decision threshold on the synthetic-validation split and
report \textbf{precision}, \textbf{recall}, and the gate's ability to separate true
physical error. We use per-axis Spearman $\rho$ between each feature and the error
to characterize which signals transfer.

\section{Experiments / Results / Discussion}

\textbf{The features see the gap (Fig.~\ref{fig:spearman}).}
\texttt{disagree\_t\_m} correlates with translation error at a remarkably stable
$\rho\approx0.59/0.56/0.59$ across synthetic/lightbox/sunlamp---the signal
transfers. The honest finding is that \texttt{disagree\_R\_deg} is essentially
\emph{dead} for rotation error ($\rho\approx0.01$/$0.08$): rotation failures are
\emph{common-mode}---symmetry flips fool both pose heads identically, so head
disagreement is structurally blind to them. Rotation signal instead comes, more
diffusely, from heatmap-peak height (\texttt{hm\_peak\_mean},
$\rho\approx{-}0.28$/${-}0.17$), heatmap entropy, and detection confidence
(\texttt{effi\_cls}).

\begin{figure}[t]
  \centering
  \includegraphics[width=0.93\linewidth]{feature_spearman.pdf}
  \caption{Per-feature Spearman $\rho$ against each error component, by domain.
  \texttt{disagree\_t\_m} tracks translation error stably across all three domains;
  \texttt{disagree\_R\_deg} is blind to rotation error (common-mode symmetry flips).}
  \label{fig:spearman}
\end{figure}

\textbf{The method matches the oracle with zero target labels
(Table~\ref{tab:auc}, Fig.~\ref{fig:roc}).} For translation the learned predictor
reaches AUC $0.841$ (lightbox) / $0.802$ (sunlamp), essentially identical to the
HIL-trained oracle ($0.846$/$0.799$) and well above the disagreement baseline
($0.800$/$0.766$) and random ($0.5$). The feature$\rightarrow$failure map is stable
across the gap, so synthetic-only training suffices. Rotation is harder but
\emph{not} hopeless: although the disagreement baseline collapses to chance
($0.51$/$0.54$), the full-feature model recovers moderate signal (AUC
$0.734$/$0.627$, again matching the oracle $0.734$/$0.631$) by combining the weak
heatmap/confidence cues above.

\begin{table}[t]
  \centering
  \small
  \caption{Failure-prediction AUC-ROC by domain and axis. The supervised method
  (synthetic-only labels) matches the HIL-trained oracle and beats both baselines;
  importance weighting (IW) is inert.}
  \label{tab:auc}
  \begin{tabular}{llccccc}
    \toprule
    Domain & Axis & Random & Disagreement & \textbf{Method} & IW & Oracle \\
    \midrule
    lightbox & $\ET$ & 0.500 & 0.800 & \textbf{0.841} & 0.839 & 0.846 \\
    lightbox & $\ER$ & 0.500 & 0.506 & \textbf{0.734} & 0.732 & 0.734 \\
    sunlamp  & $\ET$ & 0.500 & 0.766 & \textbf{0.802} & 0.790 & 0.799 \\
    sunlamp  & $\ER$ & 0.500 & 0.539 & \textbf{0.627} & 0.629 & 0.631 \\
    \bottomrule
  \end{tabular}
\end{table}

\begin{figure}[t]
  \centering
  \includegraphics[width=0.93\linewidth]{roc_curves.pdf}
  \caption{ROC curves per domain $\times$ axis. The supervised method (synthetic
  labels only) overlays the oracle; the disagreement baseline is competitive for
  translation but near-diagonal for rotation.}
  \label{fig:roc}
\end{figure}

\textbf{Importance weighting is inert---and that is the point
(Fig.~\ref{fig:iw}).} The shift is real and large: the domain classifier separates
synthetic from HIL at AUC $0.82$ (lightbox) / $0.91$ (sunlamp), and the resulting
weights collapse the effective sample size to $24\%$/$12\%$. The weights do
dramatic things, yet AUC does not move (Table~\ref{tab:auc}, IW vs.\ Method). This
confirms the covariate-shift assumption empirically: correcting $p(x)$ is
unnecessary precisely because the feature$\rightarrow$failure relationship $p(y\mid
x)$ is already stable. IW costs nothing and is the mechanism that would activate
under severe, unseen (e.g.\ on-orbit) shift.

\begin{figure}[t]
  \centering
  \includegraphics[width=0.62\linewidth]{iw_weights.pdf}
  \caption{Importance-weight distribution. Large, skewed weights (ESS collapsing to
  $12$--$24\%$) reshape the synthetic training set aggressively, yet leave AUC
  unchanged---evidence that $p(y\mid x)$ is stable across the gap.}
  \label{fig:iw}
\end{figure}

\textbf{A usable gate, not just a ranker (Table~\ref{tab:gate}).} Freezing the
decision threshold on the synthetic-validation split and applying it to HIL, the
translation gate is genuinely useful: precision $0.75$/$0.85$ at recall $\sim$0.62.
The rotation gate is marginal (precision $0.55$/$0.59$, recall $0.42$/$0.47$),
consistent with the weaker rotation signal. We note these thresholds were tuned on a
low-base-rate synthetic split; the base-rate mismatch with HIL partly explains the
modest recall, and per-domain threshold calibration is a clear lever.

\begin{table}[t]
  \centering
  \small
  \caption{Frozen-threshold gate (threshold set on synthetic val, applied to HIL).}
  \label{tab:gate}
  \begin{tabular}{llccc}
    \toprule
    Domain & Axis & Precision & Recall & Flagged rate \\
    \midrule
    lightbox & $\ET$ & 0.753 & 0.634 & 0.345 \\
    lightbox & $\ER$ & 0.550 & 0.471 & 0.235 \\
    sunlamp  & $\ET$ & 0.847 & 0.622 & 0.463 \\
    sunlamp  & $\ER$ & 0.591 & 0.423 & 0.324 \\
    \bottomrule
  \end{tabular}
\end{table}

\textbf{The gate separates real physical error (Fig.~\ref{fig:risk}).} The payoff is
in physical units: images the gate \emph{rejects} carry $\sim$3--$4.5\times$ the
true translation error of accepted ones ($0.36$ vs.\ $0.08$\,m lightbox; $0.35$ vs.\
$0.12$\,m sunlamp), and predicted failure probability rank-correlates with true
error (translation $\rho\approx0.66$/$0.63$). Crucially, this is a calibrated
\emph{risk ranking}, not an error regressor: the score is the probability of
exceeding the threshold, which is exactly what a navigation filter needs to gate or
down-weight a measurement before it corrupts the state estimate. The relationship is
weakest for sunlamp rotation ($\rho=0.25$), again tracing back to the common-mode
rotation story.

\begin{figure}[t]
  \centering
  \includegraphics[width=0.93\linewidth]{risk_vs_error.pdf}
  \caption{Predicted failure probability vs.\ true physical error, colored by the
  frozen gate decision. Rejected (red) points concentrate above the physical
  failure line; predicted risk rises with true error for translation.}
  \label{fig:risk}
\end{figure}

\textbf{Overfitting.} The strongest evidence against overfitting is that the
synthetic-only model matches the HIL-trained oracle on \emph{held-out HIL data}: a
model that had overfit synthetic idiosyncrasies would underperform the oracle, not
equal it. We additionally use $\ell_2$ regularization and a 13-dimensional feature
space against tens of thousands of training examples, making capacity-driven
overfitting unlikely.

\section{Conclusion / Future Work}

We show that the test-time reliability of a frozen, state-of-the-art spacecraft pose
estimator can be predicted from single-pass, model-internal features using a simple
logistic regression trained \emph{without any deployment-domain labels}. The
predictor matches a HIL-trained oracle for translation and recovers useful rotation
signal, functions as a deployable accept/reject gate, and isolates measurements
carrying several-fold higher physical error---directly usable by a downstream
navigation filter. Importance weighting is inert under this (real but benign) shift,
serving as free insurance for harsher regimes. Logistic regression was the
highest-value model precisely because the binding constraint is \emph{transfer},
not capacity---a finding reinforced by the oracle match. Future work: (i) a feature
that is \emph{not} common-mode to rotation symmetry flips, which is the key blocker
for rotation confidence; (ii) a constructed severe-shift test to convert IW's
``does no harm'' into a demonstrated payoff; and (iii) sequential/online arrival of
test images toward closed-loop UKF integration, where the score modulates the
measurement-noise covariance.

\section*{Contributions}
\textbf{Anjali Sreenivas}: SPNv2 inference and data collection; error-distribution
analysis; baseline experiments; writing of motivation, methods, and dataset.
\textbf{Lundeen Cahilly}: originated the confidence-prediction direction; feature
extraction from SPNv2; supervised and importance-weighted models; analysis pipeline,
gate/risk figures, and results writing.

\bibliographystyle{acl_natbib}
\bibliography{references}

\appendix
\section{Appendix}

\textbf{Preprocessing.} The \texttt{reject} feature is zero-variance on synthetic and
dropped. The \texttt{disagree\_*} features are undefined on \texttt{reject}$=1$ rows
and median-imputed from synthetic statistics. All features are standardized using
synthetic-split mean/standard deviation only, so no deployment-domain information
leaks into preprocessing.

\textbf{Calibration ablation.} Restricting the predictor to single-pass confidence
read-outs (2 features) reaches translation AUC $0.74$/$0.65$; the full 13-feature
set reaches $0.84$/$0.80$, confirming that cross-head disagreement and geometric
covariates add substantial signal beyond raw confidence scores.

\end{document}
